\documentclass[fontsize=10pt,oneside,parskip=half]{scrreprt}
\usepackage{ulille-rapport}
\usepackage{amsmath}
\usepackage{float}
\usepackage{booktabs}
\usepackage{tabularx}
\usepackage{listings}

\newcommand{\ic}[1]{{\setlength{\fboxsep}{2pt}\colorbox{BeigeSerein}{\texttt{#1}}}}

\subject{SAÉ 2.01-2.02 : Rapport POO}
\title{Versions 1 et 2}
\author{Yann Renard \\ Tom Cox \\ Alexandre Sorel}
\direction{Équipe E5}

\lstdefinestyle{bash}{
  language=bash,
  basicstyle=\ttfamily\small,
  keywordstyle=\bfseries\color{IndigoInclusif},
  commentstyle=\itshape\color{GrisStatutaire},
  stringstyle=\color{VertAvenir},
  backgroundcolor=\color{BeigeSerein},
  frame=single,
  framesep=6pt,
  rulecolor=\color{BleuHorizon},
  breaklines=true,
  breakatwhitespace=true,
  showstringspaces=false,
  tabsize=2,
  numbers=left,
  numberstyle=\tiny\color{GrisStatutaire},
  numbersep=8pt,
  xleftmargin=20pt,
  xrightmargin=4pt,
}

\lstdefinestyle{java}{
  language=Java,
  basicstyle=\ttfamily\small,
  keywordstyle=\bfseries\color{IndigoInclusif},
  commentstyle=\itshape\color{GrisStatutaire},
  stringstyle=\color{VertAvenir},
  backgroundcolor=\color{BeigeSerein},
  frame=single,
  framesep=6pt,
  rulecolor=\color{BleuHorizon},
  breaklines=true,
  breakatwhitespace=true,
  showstringspaces=false,
  tabsize=2,
  numbers=left,
  numberstyle=\tiny\color{GrisStatutaire},
  numbersep=8pt,
  xleftmargin=20pt,
  xrightmargin=4pt,
}

\begin{document}
\selectlanguage{french}

\maketitle
\pagecolor{Blanc}
{%
  \setlength{\parskip}{0pt}%
  \tableofcontents%
}

% ─────────────────────────────────────────────
\chapter{Version 1 : Un seul moyen de transport}
% ─────────────────────────────────────────────

% ──────────────────────────────
\section{Lancement et utilisation}
% ──────────────────────────────

L'application est développée avec Java~21 et Maven. Le code source est disponible
sur le dépôt GitLab de l'équipe~E5. Pour lancer l'application, il faut d'abord
compiler le projet avec Maven, puis exécuter la classe principale.

Compilation et exécution depuis la racine du projet :

\begin{lstlisting}[style=bash]
mvn clean exec:java
\end{lstlisting}

Les fichiers de données CSV doivent être placés dans \ic{src/main/resources/}.
La bibliothèque graphes (\ic{sae\_s2\_2026.jar}, \ic{jgrapht-core-1.5.1.jar},
\ic{jheaps-0.14.jar}) est référencée dans le dossier \ic{lib/} à la racine du
projet et déclarée dans le \ic{pom.xml}.

% ──────────────────────────────
\section{Diagramme UML et mécanismes objets}
% ──────────────────────────────

\noindent\textbf{Tableau 1 : Classes de la Version~1}

\begin{table}[H]
  \centering
  \begin{tabularx}{\textwidth}{llX}
    \toprule
    \textbf{Classe/Enum} & \textbf{Package} & \textbf{Rôle} \\
    \midrule
    \ic{TypeCout}    & models & Enum représentant les 3 critères : TEMPS, PRIX, CO2 \\
    \ic{Cout}        & models & Stocke les 3 valeurs de coût dans une \ic{Map<TypeCout, Double>} \\
    \ic{Ville}       & models & Implémente \ic{Lieu} (lib IUT). Représente une ville du réseau \\
    \ic{Trajet}      & models & Implémente \ic{Connexion} (lib IUT). Connexion entre deux \ic{Ville} avec un \ic{Cout} et une \ic{ModaliteTransport} \\
    \ic{Voyageur}    & models & Représente l'utilisateur : nom et \ic{TypeCout} (critère à optimiser) \\
    \ic{Plateforme}  & models & Classe centrale : contient les listes de \ic{Ville} et \ic{Trajet}, expose les fonctionnalités de chargement, filtrage et tri \\
    \bottomrule
  \end{tabularx}
  \caption{Classes de la Version~1}
\end{table}

\begin{figure}[H]
  \centering
  \includegraphics[width=0.85\textwidth]{imgrapport/diagramme_uml_v1}
  \caption{Diagramme UML — Version~1}
\end{figure}

\subsection*{Mécanismes objets mis en \oe uvre}

\noindent\textbf{Interfaces et implémentation} : Les classes \ic{Ville} et
\ic{Trajet} implémentent les interfaces \ic{Lieu} et \ic{Connexion} fournies par
la bibliothèque de l'IUT. Ce choix permet à notre code d'être directement compatible
avec la bibliothèque de graphes, qui attend des objets \ic{Lieu} et \ic{Connexion}.
C'est le principe de programmation par interface : on dépend d'une abstraction
plutôt que d'une implémentation concrète.

\noindent\textbf{Énumérations} : \ic{TypeCout} et \ic{ModaliteTransport} (fournie
par la lib) sont des enums. Ce mécanisme est approprié ici car les valeurs possibles
sont fixes et connues à l'avance. Utiliser une enum plutôt qu'une \ic{String} évite
les erreurs de frappe et rend le code plus lisible.

\noindent\textbf{Collections} : La classe \ic{Plateforme} utilise des
\ic{List<Ville>} et \ic{List<Trajet>} pour stocker le réseau. La classe \ic{Cout}
utilise une \ic{Map<TypeCout, Double>} pour associer chaque type de coût à sa valeur.
Ce choix permet d'accéder à n'importe quel critère avec \ic{getValeur(TypeCout)}.

\noindent\textbf{Encapsulation} : Tous les attributs sont privés. L'accès se fait
uniquement via des getters, ce qui protège l'état interne des objets et permet de
modifier l'implémentation sans impacter le code appelant.

% ──────────────────────────────
\section{Analyse technique des fonctionnalités}
% ──────────────────────────────

\noindent\textbf{Tableau 2 : Correspondance fonctionnalités / méthodes}

\begin{table}[H]
  \centering
  \begin{tabularx}{\textwidth}{lX}
    \toprule
    \textbf{Méthode} & \textbf{Fonctionnalité couverte (sujet)} \\
    \midrule
    \ic{chargerDepuisTableau(String[])} & Vérifier la validité des données et reconstituer le réseau \\
    \ic{filtrerParModalite(ModaliteTransport)} & Filtrer les données sur un moyen de transport spécifique \\
    \ic{getVille(String)} & Utilitaire interne pour retrouver une ville par son nom \\
    \ic{cheminExiste(...)} & Déterminer s'il existe un voyage possible entre deux villes \\
    \ic{getTrajetsTries(TypeCout)} & Afficher les meilleurs voyages ordonnés selon le critère \\
    \ic{filtrerParBorne(...)} & Exclure les alternatives qui excèdent les bornes définies \\
    \bottomrule
  \end{tabularx}
  \caption{Correspondance fonctionnalités / méthodes}
\end{table}

La méthode \ic{chargerDepuisTableau} effectue la validation des données en
vérifiant que chaque ligne contient exactement 6 colonnes et que tous les coûts
sont positifs ou nuls. En cas d'erreur, une \ic{IllegalArgumentException} est
levée (remplacée par \ic{DonneesInvalidesException} en Version~2).

La recherche des meilleurs itinéraires repose sur la bibliothèque de graphes : on
construit un \ic{MultiGrapheOrienteValue} dont les sommets sont les \ic{Ville} et
les arêtes les \ic{Trajet} valuées selon le critère choisi, puis on appelle
\ic{AlgorithmeKPCC.kpcc} pour obtenir les $k$ plus courts chemins.

% ──────────────────────────────
\section{Analyse des tests}
% ──────────────────────────────

Deux classes de test ont été développées avec JUnit~5 : \ic{CoutTest} et
\ic{PlateformeTest}.

\noindent\textbf{\ic{CoutTest}} couvre :
\begin{itemize}
  \item La création d'un \ic{Cout} avec des valeurs valides
  \item La récupération correcte de chaque valeur (PRIX, TEMPS, CO2)
  \item Le cas limite avec des valeurs à zéro
\end{itemize}

\noindent\textbf{\ic{PlateformeTest}} couvre :
\begin{itemize}
  \item Le chargement correct des données : vérification du nombre de villes (4) et de trajets (6)
  \item La détection des données invalides : ligne incomplète, coûts négatifs
  \item Le filtrage par modalité : vérification du nombre de trajets TRAIN et AVION
  \item La recherche de ville : ville existante et ville inexistante (\ic{null})
  \item La vérification d'existence de chemin : chemin existant, chemin inexistant, ville inconnue
  \item Le tri par critère : vérification que le premier élément est bien le moins coûteux
  \item Le filtrage par borne : vérification que tous les trajets retournés respectent la limite
\end{itemize}

Ces tests couvrent les cas simples, les cas limites et les cas d'erreur pour chaque
fonctionnalité principale de cette version.

% ─────────────────────────────────────────────
\chapter{Version 2 : Multi-modalité}
% ─────────────────────────────────────────────

% ──────────────────────────────
\section{Lancement et utilisation}
% ──────────────────────────────

Le lancement est identique à la Version~1. Les nouvelles fonctionnalités (import CSV,
exceptions, points d'intérêt) s'utilisent directement via la classe \ic{Plateforme}.
Les fichiers CSV de test sont disponibles dans \ic{src/test/resources/} :

\begin{itemize}
  \item \ic{exemple-reseau.csv} : données du réseau au format
        \ic{villeDépart;villeArrivée;modalité;prix;co2;durée}
  \item \ic{exemple-correspondances.csv} : coûts de correspondance au format
        \ic{ville;mode1;mode2;durée;co2;prix}
\end{itemize}

Exemple d'utilisation depuis le code :

\begin{lstlisting}[style=java]
Plateforme p = new Plateforme();
p.chargerDepuisCSV("src/test/resources/exemple-reseau.csv");
\end{lstlisting}

% ──────────────────────────────
\section{Diagramme UML et mécanismes objets}
% ──────────────────────────────

\noindent\textbf{Tableau 3 : Nouveautés de la Version~2}

\begin{table}[H]
  \centering
  \begin{tabularx}{\textwidth}{lX}
    \toprule
    \textbf{Élément ajouté} & \textbf{Justification (sujet)} \\
    \midrule
    \ic{AucunCheminException}      & Signaler l'absence de chemin via un mécanisme d'exception \\
    \ic{DonneesInvalidesException} & Signaler les problèmes de validité de données via un mécanisme d'exception \\
    \ic{chargerDepuisCSV(String)}  & Import des données issues de fichiers CSV \\
    \ic{getPointsInteret(List<Trajet>)} & N'afficher que le départ, l'arrivée et les lieux où il y a un changement de modalité \\
    \bottomrule
  \end{tabularx}
  \caption{Nouveautés de la Version~2}
\end{table}

\begin{figure}[H]
  \centering
  \includegraphics[width=0.85\textwidth]{imgrapport/diagramme_uml_v2}
  \caption{Diagramme UML — Version~2}
\end{figure}

\subsection*{Mécanismes objets mis en \oe uvre}

\noindent\textbf{Exceptions personnalisées} : Deux classes d'exception ont été créées
en étendant \ic{Exception}. \ic{AucunCheminException} est levée lorsque les villes
demandées sont inconnues ou qu'aucun trajet ne les relie. \ic{DonneesInvalidesException}
est levée lorsque les données sont mal formatées ou contiennent des valeurs invalides.
Ce mécanisme permet à l'appelant de distinguer les deux types d'erreur et de les traiter
séparément.

\noindent\textbf{Restructuration de \ic{cheminExiste}} : La méthode ne retourne plus
\ic{false} en cas d'absence de chemin : elle lève systématiquement une
\ic{AucunCheminException}. Ce choix est plus cohérent avec le principe selon lequel
une méthode nommée \ic{cheminExiste} ne devrait jamais retourner \ic{false} si
l'appelant s'attend à ce qu'un chemin existe.

\noindent\textbf{Lecture de fichier} : La méthode \ic{chargerDepuisCSV} utilise un
\ic{BufferedReader} avec try-with-resources, ce qui garantit que le fichier est
toujours fermé proprement, même en cas d'exception. Les \ic{IOException} sont
converties en \ic{DonneesInvalidesException} pour homogénéiser la gestion des erreurs.

% ──────────────────────────────
\section{Analyse technique des fonctionnalités}
% ──────────────────────────────

La méthode \ic{getPointsInteret} prend en entrée une liste de \ic{Trajet}
représentant un itinéraire complet et retourne uniquement les points significatifs :
le premier trajet, le dernier, et ceux dont la modalité diffère du trajet précédent.
Cette méthode répond à l'exigence du sujet : \emph{« n'afficher que le départ,
l'arrivée et les lieux où il y a un changement dans la modalité de transport »}.

Les exceptions personnalisées améliorent la lisibilité du code appelant : au lieu
d'attraper une \ic{IllegalArgumentException} générique, l'appelant peut traiter
spécifiquement les erreurs de données invalides et les absences de chemin.

% ──────────────────────────────
\section{Analyse des tests}
% ──────────────────────────────

Les classes de test existantes ont été mises à jour pour couvrir les nouvelles
fonctionnalités.

\noindent\textbf{\ic{PlateformeTest} — nouveaux cas de test :}
\begin{itemize}
  \item \ic{chargerDepuisCSV} : vérification du chargement depuis un fichier CSV valide
  \item \ic{cheminExiste} avec exception : vérification que \ic{AucunCheminException}
        est bien levée pour une ville inconnue et pour deux villes sans connexion directe
  \item \ic{getPointsInteret} : vérification que seuls le départ, l'arrivée et les
        changements de mode sont retournés
\end{itemize}

\end{document}
